\subsection{Třída Analyzer a režimy analýzy}
\label{subsec:implementation-analyzer}

Komunikace s jazykovým modelem je zapouzdřena ve třídě \texttt{Analyzer} v souboru \texttt{llm\_reviewer.py}.
Třída přijímá seznam objektů \texttt{EmbeddedFile} připravených v souboru \texttt{job.py}, nakonfigurovaného klienta knihovny \texttt{openai} a zvolený režim analýzy.
Knihovna \texttt{openai} poskytuje jednotné rozhraní pro všechny OpenAI-kompatibilní servery, takže stejný kód funguje jak pro lokální instanci Ollama, tak pro cloudové poskytovatele.

Centrálním vstupním bodem třídy je metoda \texttt{process}, která na základě zvoleného režimu orchestruje celý průběh analýzy.
Při volbě režimu \texttt{ZERO\_SHOT} nebo \texttt{THINKING} je provedeno jediné volání příslušné analytické metody.
Při volbě \texttt{CHAIN\_OF\_THOUGHT} jsou postupně spuštěny tři fáze -- \textbf{hrubá analýza}, \textbf{kritická revize} a \textbf{finální shrnutí} -- přičemž výstup každé fáze je předán jako kontext fázi následující.
Srovnání režimů z pohledu kvality výstupů a doby zpracování je uvedeno v \secref[sekci]{subsec:choice-strategies}.

\subsubsection*{Normalizace názvu souboru}
\label{subsubsec:implementation-analyzer-normalization}

Po dokončení analýzy je výsledek předán metodě \texttt{normalize\_issue\_filenames}.
Jazykové modely v názvech souborů občas vynechávají část cesty nebo ji mírně pozměňují.
Metoda proto každý vrácený název porovná se skutečnými vstupními cestami pomocí postupné shody přípon, prefixů a základního názvu a při jednoznačné shodě název opraví.
Tím je zajištěno, že každý nalezený problém odkazuje na platný soubor bez ohledu na to, jak jeho cestu model zapsal.

Toto opatření reaguje na konkrétní případy zjištěné během testování.
Studenti občas odevzdají soubor, ke kterému operační systém při duplicitních názvech připojí číselnou příponu (např.\ \texttt{main.c~(1)}).
Některé modely tuto příponu ve svém výstupu vynechávají, což následně vede k chybám při mapování problému na správný soubor.
V jiných případech modely nahrazují původní název generickým; například soubor \texttt{MIK0486.c} byl ve výstupu některých modelů uveden pouze jako \texttt{main.c}.

\endinput